\chapter{Background and Related Work}
\label{ch:background}

This chapter reviews the work this dissertation builds on. It covers the class of model under study, the
difficulties that make animal detection and tracking hard in the first place, how tracking is scored, the
literature on estimating a model's performance without labels, what is known about distribution shift in
camera-trap vision, the representations that make a distance between species meaningful, and the datasets
involved. It closes by stating the gap that remains.

\section{Promptable and open-vocabulary segmentation and tracking}
\label{sec:bg-promptable}
For most of the last decade, detection and segmentation models were trained against a fixed list of
categories. Such a model cannot be asked about anything outside that list. Extending it to a new category
means collecting labelled examples and training again. In camera-trap work the standard tool of this kind is
MegaDetector~\cite{megadetector}, a detector trained on millions of camera-trap images to localise animals,
people and vehicles, and widely deployed as a first-pass filter over survey footage.

The \emph{Segment Anything} line broke the fixed-list assumption. It replaced the label set with a prompt: a
user indicates what should be segmented, and the model returns a mask for it. SAM\,2~\cite{sam2} carried the
formulation from images to video, propagating a mask across frames from a point or box placed in one of them.

SAM\,3~\cite{sam3} changes what the prompt may be. Instead of a point or a box it accepts a short noun
phrase, and returns every instance matching that phrase, segmented and tracked across the clip. The category
need never have been seen in training. This is a zero-shot, open-vocabulary detect--segment--track
capability, and SAM\,3 is the current state of the art for it.

Other models reach open-vocabulary video understanding by different routes. GLEE~\cite{glee} trains a single
unified model to detect, segment and track objects at scale, driven by a text query. Florence-2~\cite{florence2}
takes a caption-style approach, using one sequence-to-sequence model across a range of vision tasks and
emitting regions in response to a task prompt, though without a confidence score attached to them. These are
the main alternative designs in the same space.

\section{What makes animals hard to detect and track}
\label{sec:bg-challenges}
Wild footage is an unusually hostile input, and the difficulties are worth setting out because they are what
any tracker deployed on it must survive.

The first is that animals are often hard to see at all. Many are camouflaged against the habitat they live
in, and the resulting low contrast between object and background defeats methods that assume a visually
salient foreground. This is a recognised problem in its own right: camouflaged object detection is studied as
a separate task with dedicated benchmarks, precisely because standard detectors underperform on
it~\cite{cod}.

The second is variability of appearance. A camera trap is triggered by whatever passes, so an animal may fill
the frame in one clip and occupy a few dozen pixels in the next. It may face the lens or present only a
flank, and vegetation frequently occludes part of it. A bounding box is a poor description of an animal in
such conditions, which is one reason segmentation --- labelling the pixels belonging to the animal --- is the
more informative output.

The third is image quality. Much wildlife footage is captured at night under infrared illumination, where
colour information is absent and contrast is poor. Infrared and thermal imaging are established wildlife
monitoring modalities with their own detection characteristics, and their properties differ enough from
daylight imagery that observation strategies are designed around them~\cite{thermalwildlife}. Motion blur is
common, since animals are typically moving when they trigger the camera.

The fourth is context. A camera trap is fixed, so the same scene recurs in every clip from that site, and a
model can come to rely on the scene rather than on the animal. This failure mode has a literature of its own,
reviewed in Section~\ref{sec:bg-cameratrap}.

The fifth applies to tracking rather than detection. Keeping an identity attached to the right animal is a
distinct problem from finding it. Animals leave the frame and re-enter it, pass behind vegetation, and in
groups may be near-indistinguishable from one another. Identity errors are treated as a first-class failure
mode in multi-object tracking benchmarks, which annotate visibility explicitly and report identity switches
separately from misses~\cite{motchallenge}. Recovering the identity of a specific individual across sightings
is hard enough in wildlife imagery to constitute its own research area~\cite{animalreid}.

\section{Evaluating multi-object tracking}
\label{sec:bg-hota}
Multi-object tracking is scored by comparing predicted tracks against annotated ones. The difficulty is that
two very different failures --- missing the object, and swapping its identity --- both reduce a single
accuracy figure, and earlier metrics combined them in ways that were hard to interpret.

Higher Order Tracking Accuracy (HOTA)~\cite{hota} was introduced to separate them. At each localisation
threshold $\alpha$ it forms a detection term (\texttt{DetA}: were the right objects found?) and an
association term (\texttt{AssA}: were identities kept consistent over time?), then combines them as a
geometric mean:
\begin{equation}
  \mathrm{HOTA}_\alpha=\sqrt{\mathrm{DetA}_\alpha\cdot\mathrm{AssA}_\alpha},
  \qquad
  \mathrm{pHOTA}=\frac{1}{|\mathcal{A}|}\sum_{\alpha\in\mathcal{A}}\mathrm{HOTA}_\alpha ,
  \label{eq:phota}
\end{equation}
averaged over the threshold set $\mathcal{A}=\{0.05,0.10,\dots,0.95\}$. Because the two terms are formed
separately before being combined, a HOTA score can be read either as a single number or as two.

SA-FARI scores promptable tracking with \texttt{pHOTA}, a promptable variant that decomposes identically into
\texttt{pDetA} and \texttt{pAssA}.

\section{Predicting model performance without labels}
\label{sec:bg-predictable}
A now-mature literature shows that a model's performance under distribution shift can be anticipated without
labelling the target data.

\citet{aotl} observe that out-of-distribution (OOD) accuracy is strongly linearly correlated with
in-distribution accuracy, across many models and many shifts. \citet{agl} show that the agreement between a
pair of models tracks OOD accuracy closely enough to estimate it from unlabelled data alone. A parallel
strand reads the model's own statistics on the unlabelled target. Average thresholded confidence~\cite{atc}
calibrates a confidence threshold on source data and reads off target accuracy as the mass above it, while
\citet{autoeval} regress classifier accuracy directly on a distribution-shift feature --- the Fr\'echet
distance between train and test feature sets.

The idea has only recently reached structured outputs. \citet{pcr} estimate object-detection performance
without ground truth, from the spatial consistency and confidence reliability of pre-NMS boxes. Reverse
classification accuracy~\cite{rca} predicts per-image segmentation quality by training a reverse model
against a labelled reference.

Two properties characterise essentially all of this work. First, the target is an image classifier scored by
a single scalar, or at most a detector or segmenter operating on a static image. None of it concerns identity
maintained over time, and none decomposes the estimate into separate components. Second, every one of these
methods runs the model on the target and inspects what comes out --- its predictions, its confidences, or its
agreement with another model. Each estimates the score of a model that has already been executed.

\section{Distribution shift in camera-trap and wildlife vision}
\label{sec:bg-cameratrap}
That camera-trap models degrade at new locations is well established. \emph{Terra Incognita}~\cite{terra}
showed that recognisers excellent at trained camera sites fail sharply at unseen ones. WILDS~\cite{wilds}
consolidated this, via iWildCam, as a canonical real-world shift. PanAf-FGBG~\cite{panaf} demonstrated that
behaviour-correlated backgrounds causally drive the drop, and proposed a latent-space mitigation.

The claim that scene context, rather than the object, drives such failures has a longer pedigree in general
vision. \citet{xiao} isolate foreground from background on ImageNet, in the ImageNet-9 ``Background
Challenge'', and show that a classifier scores well above chance from the background \emph{alone}. An
adversarially chosen background flips a correctly classified foreground up to $87.5\%$ of the time.
Background is genuine signal, not noise. \citet{elephant} make the parallel point for detectors:
transplanting an object into an unfamiliar scene suppresses its detection, flips its predicted identity as it
moves, and disturbs other detections in the image, exposing a reliance on context and co-occurrence.

\citet{xiao} also report that more accurate and more robust models rely on background \emph{less}, which
suggests the effect is not fixed but a property of how strong the model is.

What this literature shares is its direction of travel. It measures degradation after the fact, or intervenes
on the input to mitigate it. It does not forecast per-site or per-species reliability in advance. It also
studies classifiers and behaviour recognisers rather than promptable trackers, scored by a single accuracy
with no detection/association split.

\section{Representations of species and scenes}
\label{sec:bg-representations}
Measuring how far one species sits from another, or one scene from another, requires a representation in
which such a distance means something.

Vision foundation models turn out to carry that structure. BioCLIP~\cite{bioclip} shows that a model trained
on a large biological image corpus encodes taxonomic structure across the tree of life, its feature space
reflecting the hierarchy that runs from kingdom down to species. Appearance and phylogeny are correlated, so
a representation that separates species visually will tend also to order them taxonomically.

General-purpose encoders offer the same service without biological supervision. DINOv2~\cite{dinov2} learns
visual features self-supervised, and they transfer across domains without fine-tuning. CLIP~\cite{clip}
aligns images with natural-language descriptions, producing a space in which images and text share
coordinates. Both are widely used as off-the-shelf feature extractors.

\section{Datasets}
\label{sec:bg-safari}
Two datasets are central to this work.

\paragraph{SA-FARI.} SA-FARI~\cite{safari} is a large open multi-animal camera-trap tracking dataset,
released alongside SAM\,3. It contains $11{,}609$ densely annotated videos spanning $99$ species, recorded at
$741$ locations across four continents. Every species carries a seven-level biological taxonomy running from
kingdom to species, and every video carries a creation timestamp. Annotations are per-frame masks with
persistent identities, and the benchmark ships an official evaluator, \texttt{VEval}.

Two features set it apart. The taxonomy makes the relationship between two species measurable rather than a
matter of judgement, and the location and timestamp metadata do the same for place and era. The benchmark
also includes abundant \emph{hard negatives}: prompts naming a species that is absent from the clip, which a
correct run should answer with nothing. These make it possible to study false alarms alongside misses.

SA-FARI ships a train/test partition drawn along locations, so the test camera sites are unseen while most
species occur on both sides of it.

\paragraph{BURST.} BURST~\cite{burst} is a video benchmark for object recognition, segmentation and tracking,
built on the TAO~\cite{tao} video corpus. It annotates $482$ categories from the LVIS~\cite{lvis} vocabulary
with per-frame masks and persistent track identities, and a portion of that vocabulary is animal categories.
Its footage comes from internet, driving and film sources rather than camera traps, and it carries neither
location nor timestamp metadata. It is mask-native and considerably more varied in subject matter than
SA-FARI, which makes it a useful independent point of comparison.

\section{Summary and gap}
\label{sec:bg-gap}
Predictable OOD generalisation is established for classifiers scored by a scalar~\cite{aotl,agl,atc,autoeval}
and, very recently, for static detectors and segmenters~\cite{pcr,rca}. Camera-trap vision has thoroughly
documented, but not forecast, location-driven degradation~\cite{terra,wilds,panaf}. Most tellingly, a recent
unified study of camera-trap recognition over time~\cite{jeon} explicitly names ``predicting when zero-shot
models will succeed at a new site'' as an open question, and stops at naming it.

No prior work predicts, \emph{before running the model} and from label-free distances, the zero-shot transfer
of a promptable video tracker. None separates that prediction into detection and association. That is the gap
this dissertation addresses.
